\section{Methods and Approach}

We ask two questions: does combining flow and authentication logs in a graph improve detection over single-source methods, and how does this approach compare to classical anomaly detection on tabular features?

\subsection{Design Rationale}

Our pipeline is built around how lateral movement actually works. Attackers rely on TCP-based services --- SMB (port 445), SSH (port 22), RDP (port 3389) --- because these protocols give them reliable, ordered delivery for remote command execution and file transfer. During reconnaissance, a compromised host connects to many destinations in quick succession (high fan-out), then narrows to specific targets once a path opens up. Attacker timing also differs from normal users: lateral movement tends to involve deliberate, spaced-out connections rather than the rapid-fire patterns of routine traffic.

These patterns shape our design. We include protocol and port features in edge scoring because lateral movement is inherently service-specific. We compute fan-out and degree features at the node level to catch the scanning-to-targeting transition. We extract temporal features (inter-arrival times, burst scores, active duration) to separate deliberate attacker behavior from routine traffic.

\subsection{Graph Construction}

We model the network as a directed multigraph using the igraph library. Nodes are entities (computers and users); edges are interactions (authentication events and network flows). The graph is built incrementally by a \texttt{StreamingGraphBuilder} that reads log files line-by-line through time windows derived from the red-team ground truth. Events outside all windows are discarded immediately, keeping memory proportional to the active window set rather than the full dataset. The pipeline processes over 135 million events in approximately 4 GB of RAM.

For authentication events, the builder creates a computer-to-computer edge (source computer to destination computer) and optionally a user-to-user edge. Each edge stores authentication type, logon type, orientation, success status, and timestamp. When multiple events share the same source-destination pair, the edge weight increments and the first and last timestamps are preserved --- duplicates become weight, not separate edges. For flow events, a computer-to-computer edge stores protocol, source and destination ports, packet count, byte count, duration, and timestamp, with the same deduplication. Node types are inferred heuristically: node names ending with \texttt{\$} or containing \texttt{\$} before \texttt{@} are machine accounts; others are users. After all events are processed, the edge map is materialized into an igraph object.

Three graph variants are constructed per run: \emph{flow-only} (approximately 11,586 nodes, 131,562 edges), \emph{auth-only} (approximately 90,783 nodes, 409,100 edges), and \emph{combined} (approximately 91,589 nodes, 512,529 edges). Each graph is freed after scoring and evaluation.

\subsection{Feature Extraction}

We extract structural, temporal, and statistical features at the node, edge, and graph levels. All features are computed independently for each graph variant.

\textbf{Node features.}
For each node, we compute in-degree, out-degree, total degree, and fan-out ratio (out-degree divided by total degree). The fan-out ratio targets a specific lateral movement pattern: compromised hosts performing reconnaissance connect to many distinct destinations, while normal servers mostly receive inbound connections. Betweenness centrality is computed exactly for graphs under 5,000 nodes; for larger graphs, we use an approximation with a path length cutoff of 3 hops, trading some accuracy for tractability.

Temporal features capture timing patterns of outgoing connections. For each node, we compute mean and standard deviation of inter-arrival times, a burst score (the maximum fraction of outgoing edges falling within any 10\% sub-window of the node's active span), and total active duration. The intuition is that reconnaissance produces bursty, rapid-fire connections while persistent lateral movement shows longer, more regular intervals.

\textbf{Edge features.}
For each edge, we compute edge rarity (the reciprocal of edge weight), source out-degree, destination in-degree, and binary indicators for authentication signals: whether the authentication uses NTLM, whether the logon type is Network (remote access), and whether authentication succeeded. NTLM is particularly relevant because it is uncommon in normal traffic but frequently appears in pass-the-hash attacks.

For flow edges, we compute an unusual destination port indicator flagging ports associated with lateral movement (SSH-22, Telnet-23, SMB-445, RDP-3389, WinRM-5985/5986) or ephemeral ports above 49,152; protocol rarity (one minus the ratio of edges using that protocol to total flow edges); byte-per-packet ratio as a percentile rank; and duration z-score across all edges. Self-loops and user-to-user edges are excluded from scoring.

\textbf{Graph features.}
We compute density, average clustering coefficient, connected component count, node count, and edge count. These give a snapshot of each graph variant's structural complexity.

\subsection{Anomaly Scoring}

\subsubsection{Edge Scoring}
Edges are scored using a weighted combination of features, with the formula depending on edge type. The current weights are initial values chosen based on domain knowledge; we plan to tune them systematically in future work.

Authentication edges:
\begin{equation}
s_{\text{auth}} = 0.4 \cdot \text{is\_ntlm} + 0.3 \cdot \text{is\_network\_logon} + 0.3 \cdot \text{rarity\_rank}
\end{equation}
NTLM receives the highest weight because pass-the-hash attacks are a common lateral movement technique. Network logon type indicates remote access rather than interactive sessions. Edge rarity (percentile-ranked) identifies uncommon source-destination pairs.

Flow edges:
\begin{equation}
s_{\text{flow}} = 0.4 \cdot \text{rarity\_rank} + 0.3 \cdot \text{is\_unusual\_dst\_port} + 0.3 \cdot \text{protocol\_rarity\_rank}
\end{equation}
Rarity is weighted highest because uncommon connections are the strongest signal available from flow data. The unusual port flag encodes knowledge of lateral movement services, and protocol rarity measures how much a given protocol deviates from the normal distribution.

For the combined method, authentication edges receive a $1.5\times$ multiplier (\texttt{auth\_weight\_multiplier}) because auth signals are more informative but far fewer than flow edges. Without this multiplier, auth scores would be swamped. Self-loops and user-to-user edges are scored as 0.

\subsubsection{Path Enumeration and Scoring}
To capture multi-step lateral movement, we enumerate paths through the graph. Starting from each node, breadth-first search explores up to 4 hops, following only the top 10 outgoing edges per node ranked by score. Enumeration runs in parallel across 12 CPU workers.

Each path is scored as the mean of three aggregates of its edge scores: geometric mean (flags paths where every edge is somewhat suspicious), maximum edge score (flags paths with at least one highly anomalous hop), and arithmetic mean. The top 50 paths are retained. Edge scores then get a boost of 10\% of any traversing path's score (\texttt{path\_boost\_factor} $= 0.1$), feeding path-level signals back into edge-level detection.

\subsection{Threshold Optimization and Detection}

Anomalous edges are identified by thresholding final edge scores. Rather than fixing a percentile, we sweep across $[90, 95, 97, 99, 99.5, 99.9]$ and select the threshold maximizing F1 against known red-team pairs. This search operates in pair-space: after thresholding, anomalous source-destination pairs are compared against ground-truth pairs to compute recall, precision, and F1. We acknowledge that optimizing the threshold on the same labeled data used for evaluation is a limitation --- the selected thresholds (95th percentile for auth-only, 97th for flow-only and combined) may not generalize to unseen attack patterns.
